\documentclass[11pt,a4paper]{article}

% Packages
\usepackage[utf8]{inputenc}
\usepackage[T1]{fontenc}
\usepackage{geometry}
\usepackage{graphicx}
\usepackage{hyperref}
\usepackage{booktabs}
\usepackage{enumitem}
\usepackage{float}
\usepackage{fancyhdr}
\usepackage{parskip}
\usepackage{titlesec}

% Compact geometry for 5-page limit
\geometry{left=2cm, right=2cm, top=2cm, bottom=2cm}

% Compact section spacing
\titlespacing*{\section}{0pt}{1.5ex plus 0.5ex minus 0.2ex}{1ex plus 0.2ex}
\titlespacing*{\subsection}{0pt}{1ex plus 0.3ex minus 0.1ex}{0.5ex plus 0.1ex}

% Header
\pagestyle{fancy}
\fancyhf{}
\rhead{COMP3011 Coursework 1}
\lhead{Technical Report}
\rfoot{\thepage}

% Hyperlink setup
\hypersetup{colorlinks=true, linkcolor=blue, urlcolor=blue}

\title{\textbf{Smart Productivity Analytics Platform API}\\[0.3cm]\large Technical Report -- COMP3011 Web Services and Web Data}
\author{Student ID: [Insert Student ID]}
\date{March 2026}

\begin{document}

\maketitle
\thispagestyle{fancy}

% =============================================================================
\section{Introduction and Project Overview}
% =============================================================================

This report documents the design, implementation, and evaluation of a RESTful Web API for a Smart Productivity Analytics Platform developed for COMP3011 coursework. The system provides task management, habit tracking, Pomodoro timer functionality, and productivity analytics through a well-documented API with AI agent integration capabilities.

The platform was designed to address the need for productivity tools that integrate with modern AI assistants while providing comprehensive data management and analytics. The implementation demonstrates proficiency in web service development, database design, authentication systems, and software testing methodologies.

% =============================================================================
\section{Technology Stack and Justification}
% =============================================================================

The technology stack was selected based on maturity, documentation quality, and suitability for rapid API development:

\begin{table}[H]
\centering
\small
\begin{tabular}{@{}llp{7cm}@{}}
\toprule
\textbf{Component} & \textbf{Choice} & \textbf{Justification} \\
\midrule
Framework & Django 5.2 & Mature, batteries-included framework with excellent ORM and admin interface \\
API Layer & Django REST Framework & Industry-standard for building RESTful APIs with serialization, authentication, and pagination \\
Authentication & Simple JWT & Stateless authentication enabling horizontal scaling; token rotation enhances security \\
Documentation & drf-spectacular & OpenAPI 3.0 schema generation with interactive Swagger UI \\
Database & SQLite/PostgreSQL & SQLite for development portability; PostgreSQL-ready for production \\
Testing & pytest-django & Flexible fixtures, parallel execution, and excellent Django integration \\
\bottomrule
\end{tabular}
\caption{Technology Stack with Justification}
\end{table}

The decision to use Django over alternatives such as Flask or FastAPI was driven by its comprehensive feature set, including built-in admin interface, robust ORM, and extensive middleware ecosystem. Django REST Framework was chosen for its mature serialization system and viewset patterns that reduce boilerplate code.

% =============================================================================
\section{API Design and Architecture}
% =============================================================================

\subsection{Architectural Approach}

The system follows a layered architecture with clear separation of concerns: URL routing directs requests to ViewSets, which delegate data transformation to Serializers, while Models handle database interactions through Django's ORM. This separation facilitates unit testing and code maintainability.

API versioning is implemented via URL prefixes (\texttt{/api/v1/}) to ensure backward compatibility for future iterations.

\subsection{RESTful Endpoint Design}

The API implements resource-oriented URLs following REST conventions. Key endpoint groups include:

\begin{itemize}[noitemsep]
    \item \textbf{Authentication}: Registration, JWT token management, profile operations
    \item \textbf{Task Management}: Full CRUD with filtering, bulk operations, and status transitions
    \item \textbf{Productivity}: Habit tracking with streak calculation, Pomodoro session management
    \item \textbf{Analytics}: Dashboard statistics and productivity trend analysis
    \item \textbf{MCP Integration}: Model Context Protocol endpoints for AI agent interoperability
\end{itemize}

HTTP methods are used appropriately: GET for retrieval, POST for creation, PUT/PATCH for updates, and DELETE for removal. Custom actions (e.g., \texttt{/tasks/\{id\}/complete/}) use POST as they trigger state changes.

\subsection{MCP Integration}

A novel feature of this implementation is Model Context Protocol (MCP) compatibility, enabling AI agents to discover and execute API operations programmatically. The \texttt{/mcp/capabilities/} endpoint exposes system metadata, while \texttt{/mcp/tools/} provides JSON Schema definitions for available operations. This design demonstrates forward-thinking integration with emerging AI assistant ecosystems.

% =============================================================================
\section{Data Model Design}
% =============================================================================

The database schema comprises seven interconnected entities designed to support comprehensive productivity tracking:

\begin{itemize}[noitemsep]
    \item \textbf{User}: Extended Django AbstractUser with productivity preferences (timezone, work hours, daily goals)
    \item \textbf{Task}: Core entity with status workflow, five-level priority system, energy level tagging, and subtask support via self-referential foreign key
    \item \textbf{Category/Tag}: Hierarchical categorization (two-level nesting) and flexible tagging with many-to-many relationships
    \item \textbf{Habit}: Frequency-based tracking with automatic streak calculation logic
    \item \textbf{PomodoroSession}: Focus session tracking with interruption counting and task association
    \item \textbf{ProductivitySnapshot}: Daily metrics for trend analysis
\end{itemize}

UUID primary keys were chosen over auto-incrementing integers to prevent enumeration attacks and support future distributed database scenarios. Database indexes on frequently-queried field combinations (user+status, user+due\_date) optimize query performance.

% =============================================================================
\section{Authentication and Security}
% =============================================================================

Security was implemented through multiple layers:

\textbf{JWT Authentication}: Access tokens expire after 60 minutes, minimizing exposure from token theft. Refresh tokens (7-day lifetime) enable seamless re-authentication. Token rotation and blacklisting prevent reuse of compromised refresh tokens.

\textbf{Password Security}: Django's built-in validators enforce minimum length, prevent common passwords, and block user-attribute similarity.

\textbf{Rate Limiting}: Anonymous requests are limited to 100/hour; authenticated users receive 1000/hour, protecting against abuse while maintaining usability.

\textbf{Data Isolation}: All querysets are filtered by the authenticated user, ensuring users cannot access others' data even with valid object IDs.

% =============================================================================
\section{Testing Strategy}
% =============================================================================

The project employs a comprehensive testing approach using pytest-django:

\textbf{Test Organization}: Tests are separated by domain---\texttt{test\_auth.py} covers authentication flows, \texttt{test\_tasks.py} validates CRUD operations and filtering, and \texttt{test\_mcp.py} verifies AI integration endpoints.

\textbf{Fixture-Based Setup}: Reusable fixtures in \texttt{conftest.py} provide authenticated clients, sample users, and test data, reducing test code duplication.

\textbf{Coverage}: The test suite includes 57 tests covering positive paths, error conditions, and edge cases. Authentication tests verify both successful flows and security boundaries (e.g., unauthenticated access returns 401).

\textbf{Integration Testing}: Tests execute full request/response cycles through Django's test client, validating serialization, database operations, and response formatting.

% =============================================================================
\section{Documentation}
% =============================================================================

API documentation is generated automatically from code using drf-spectacular, producing OpenAPI 3.0 schemas. Interactive documentation is available via Swagger UI at \texttt{/api/v1/docs/swagger/} and ReDoc at \texttt{/api/v1/docs/redoc/}.

Schema decorators (\texttt{@extend\_schema}) provide endpoint-specific descriptions, request/response examples, and parameter documentation. This approach ensures documentation remains synchronized with implementation.

% =============================================================================
\section{Use of Generative AI}
% =============================================================================

In accordance with coursework guidelines, this section declares the use of Generative AI tools during development:

\textbf{Tools Used}: GitHub Copilot (Claude) was employed throughout the development process.

\textbf{Applications}:
\begin{itemize}[noitemsep]
    \item \textbf{Code Generation}: Initial boilerplate for models, serializers, and viewsets was generated with AI assistance, then reviewed and refined
    \item \textbf{Debugging}: AI was used to diagnose issues such as drf-spectacular enum naming collisions and test fixture configuration
    \item \textbf{Documentation}: Technical report structure and content were developed with AI assistance
    \item \textbf{Test Development}: Test cases were co-developed with AI to ensure comprehensive coverage
\end{itemize}

\textbf{Critical Evaluation}: While AI assistance accelerated development, all generated code was reviewed for correctness, security implications, and adherence to project requirements. AI suggestions for authentication were particularly scrutinized to ensure proper security implementation.

% =============================================================================
\section{Conclusion}
% =============================================================================

This project delivers a production-ready RESTful API demonstrating competency in web service development. Key achievements include:

\begin{itemize}[noitemsep]
    \item Complete CRUD operations with advanced filtering and bulk operations
    \item Robust JWT authentication with security best practices
    \item Novel MCP integration for AI agent compatibility
    \item Comprehensive automated testing (57 tests)
    \item Interactive OpenAPI documentation
\end{itemize}

The implementation addresses assessment criteria across all grade bands, with particular strengths in architectural design, security implementation, and innovative AI integration features.

% =============================================================================
\section*{References}
% =============================================================================

\begin{enumerate}[noitemsep, label={[\arabic*]}]
    \item Django Documentation. \url{https://docs.djangoproject.com/}
    \item Django REST Framework. \url{https://www.django-rest-framework.org/}
    \item JSON Web Tokens Introduction. \url{https://jwt.io/introduction}
    \item OpenAPI Specification. \url{https://swagger.io/specification/}
    \item Model Context Protocol. \url{https://modelcontextprotocol.io/}
\end{enumerate}

\end{document}
